\documentclass[11pt]{article}
\usepackage[margin=0.85in]{geometry}
\usepackage{amsmath,amssymb}
\usepackage{booktabs}
\usepackage{enumitem}
\usepackage{xcolor}
\usepackage{hyperref}
\usepackage{array}
\usepackage{tabularx}

\hypersetup{colorlinks=true,linkcolor=blue,urlcolor=blue}

\newcommand{\softmax}{\operatorname{softmax}}
\newcommand{\softplus}{\operatorname{softplus}}
\newcommand{\1}{\mathbf{1}}

\title{\textbf{Intuition Behind the TP2-MV-ALiBi Transformer}\\
\large A readable explanation for the SPY call TP2 project}
\author{SPY American Call TP2 Project}
\date{\today}

\begin{document}
\maketitle

\section{Purpose of This Note}

This note explains the transformer used in the TP2 project in plain language.
The goal is not only to state the equation, but to explain what the model is
trying to learn and why the architecture was chosen.

The key point is:

\begin{center}
\textbf{The transformer is not detecting TP2 violations. The detector already
finds the violations. The transformer decides whether a detected violation is
worth trading.}
\end{center}

For each TP2 violation episode, the model asks:

\begin{center}
\textit{Given what this violation looks like today, and given what happened
to similar violations in the past, should we trade T1, trade K2, or skip?}
\end{center}

This distinction matters. A TP2 violation is a mathematical signal, but not
every mathematical violation is economically useful after bid/ask spreads,
liquidity, market regime changes, and payoff asymmetry. The machine learning
problem is therefore a conditional selection problem.

\section{Where the Transformer Sits in the Pipeline}

The pipeline before the transformer is:

\begin{enumerate}[leftmargin=*]
  \item Clean the SPY call option data.
  \item Construct option-price rectangles across two strikes and two maturities.
  \item Detect conservative bid/ask TP2 violations.
  \item Track each violation as an episode.
  \item Build strategy labels for T1 and K2 under the corrected payoff rule.
  \item Build point-in-time features available on the violation date.
\end{enumerate}

Only after these steps do we train the transformer. The model never sees future
P\&L, future exit date, test-year labels, or the year as a direct feature.

\section{What Is a Token?}

The transformer reads \textbf{tokens}. In our project, one token is one
episode-strategy candidate:

\begin{center}
\textbf{one TP2 episode + one candidate strategy, either T1 or K2.}
\end{center}

So if a TP2 violation episode is found on a certain date, we create two model
rows:

\[
  (\text{episode } i, \text{T1}), \qquad
  (\text{episode } i, \text{K2}).
\]

Each token contains only entry-time information, such as:

\begin{itemize}[leftmargin=*]
  \item TP2 severity: how strong the inequality violation is.
  \item Moneyness and distance from the forward/underlying.
  \item Time to expiration for the short and long maturity legs.
  \item Bid/ask spread, volume, and open interest.
  \item Short-D-leg information, because both T1 and K2 short the shorter-dated
        D leg.
  \item Lifecycle pressure, such as whether the violation is fresh or persistent.
  \item Strategy identity: whether the token represents T1 or K2.
\end{itemize}

The model output for each token is a probability:

\[
  p_{i,r} =
  \Pr(\text{strategy } r \text{ is profitable for episode } i
  \mid \text{entry-time information}).
\]

Then a validation-tuned rule chooses:

\[
  \text{trade T1, trade K2, or skip.}
\]

\section{Why Not Just Use a Normal MLP or LightGBM?}

An MLP or LightGBM can look at today's violation features. That is useful, but
it does not naturally ask:

\begin{center}
\textit{Have we seen similar TP2 episodes before, and what happened after them?}
\end{center}

The transformer is useful because it can compare today's episode with a short
history of past TP2 episodes. This is important because option-market behavior
is not independent day by day. Some periods have persistent liquidity stress,
some have strong trend regimes, and some have unusually profitable or
unprofitable payoff behavior.

The transformer is therefore a memory-based model:

\begin{center}
\textbf{current TP2 episode} + \textbf{similar historical TP2 episodes}
\(\rightarrow\) \textbf{trade probability}.
\end{center}

\section{The Attention Equation in Plain Language}

The core attention score is:

\[
  a_{ij}^{(h)}
  =
  \frac{(Q_i^{(h)})^\top K_j^{(h)}}{\sqrt{d_h}}
  -
  m_h^\top d_{ij}
  +
  \rho_h^{\mathrm{same}}\1\{r_i=r_j\}
  +
  \rho_h^{\mathrm{cross}}\1\{r_i\neq r_j\}.
\]

This equation decides how much current episode \(i\) should pay attention to
past episode \(j\), in attention head \(h\).

\subsection{Part 1: Feature Similarity}

\[
  \frac{(Q_i^{(h)})^\top K_j^{(h)}}{\sqrt{d_h}}
\]

This is the standard transformer similarity term.

Intuition:

\begin{itemize}[leftmargin=*]
  \item \(Q_i\) means: what the current episode is looking for.
  \item \(K_j\) means: what the past episode represents.
  \item Their dot product is large when the current and past episodes look
        similar in the learned feature space.
\end{itemize}

If today's violation resembles a past violation that was profitable, the model
can use that historical information. If it resembles a past violation that
lost money, the model can also learn caution.

\subsection{Part 2: Multivariate ALiBi Penalty}

\[
  -m_h^\top d_{ij}
\]

This is the main TP2-specific modification.

Standard attention only learns similarity. But in financial data, two episodes
can look similar numerically while still being meaningfully different. For
example:

\begin{itemize}[leftmargin=*]
  \item one happened yesterday, another happened five years ago;
  \item one was near-the-money, another was deep in-the-money;
  \item one was liquid, another had wide spreads;
  \item one was a fresh violation, another was an old persistent violation.
\end{itemize}

The distance vector \(d_{ij}\) measures these differences. In our model it
contains five domain distances:

\[
  d_{ij}
  =
  [
  \text{time distance},
  \text{surface distance},
  \text{short-D-leg distance},
  \text{liquidity distance},
  \text{lifecycle distance}
  ].
\]

The learned vector \(m_h\) says how much each attention head cares about each
kind of distance. If a head learns that liquidity mismatch is dangerous, it
will put a larger penalty on liquidity distance. If another head cares more
about maturity structure, it can penalize surface or short-D-leg distance more.

The negative sign is important:

\[
  \text{larger distance} \Rightarrow \text{lower attention}.
\]

So the model is encouraged to use past episodes that are not only similar in
features, but similar in financially meaningful ways.

\subsection{Part 3: Same-Strategy and Cross-Strategy Bonuses}

\[
  \rho_h^{\mathrm{same}}\1\{r_i=r_j\}
  +
  \rho_h^{\mathrm{cross}}\1\{r_i\neq r_j\}
\]

This tells the transformer whether it should compare T1 mostly with past T1
examples, or whether past K2 examples are also informative.

For example:

\begin{itemize}[leftmargin=*]
  \item Some heads may learn that T1 should mostly attend to T1 history.
  \item Other heads may learn that K2 history helps explain T1 outcomes because
        both strategies share the short D leg.
\end{itemize}

This is useful because the two viable strategies are related but not identical.

\section{From Attention Scores to a Decision}

After computing scores \(a_{ij}^{(h)}\), the model applies softmax:

\[
  \alpha_{ij}^{(h)}
  =
  \softmax_j(a_{ij}^{(h)}).
\]

The values \(\alpha_{ij}^{(h)}\) are attention weights. They sum to one across
the historical episodes. A high weight means:

\begin{center}
\textit{this past episode is important for understanding the current episode.}
\end{center}

The transformer then forms a historical context vector:

\[
  c_i^{(h)} = \sum_j \alpha_{ij}^{(h)} V_j^{(h)}.
\]

In plain language:

\begin{center}
\textbf{The model builds a weighted memory of the most relevant past TP2
episodes.}
\end{center}

Finally, the current episode representation and the historical memory are
combined and passed through a classifier:

\[
  \Pr(y_{i,r}=1)
  =
  \sigma(\mathrm{MLP}([h_i,c_i])).
\]

Here \(y_{i,r}=1\) means strategy \(r\) was profitable for episode \(i\) under
the chosen payoff label.

\section{Why This Architecture Makes Sense for TP2}

The TP2 opportunity is not just a static row-wise pattern. It is a market
structure pattern. The same violation severity can mean different things under
different conditions:

\begin{itemize}[leftmargin=*]
  \item In a liquid, stable market, a violation may revert or produce clean
        payoff value.
  \item In a stressed market, wide spreads and regime movement may dominate
        the theoretical signal.
  \item A fresh violation may be more informative than a violation that has
        persisted for many days.
  \item A near-expiration short-D-leg setup behaves differently from a longer
        maturity setup.
\end{itemize}

The transformer is designed to learn these conditional patterns from historical
episodes.

\section{What the Model Is Not Doing}

It is important to be precise:

\begin{itemize}[leftmargin=*]
  \item It is not proving arbitrage.
  \item It is not guaranteeing executable live trading profits.
  \item It is not using future information at entry time.
  \item It is not using LightGBM predictions as transformer inputs.
  \item It is not selecting the best strategy after seeing the test year.
\end{itemize}

The model is a conditional selector trained under rolling out-of-sample splits.
It tests whether TP2 violations contain learnable structure that can improve
trade selection.

\section{Why 2022 Is Difficult}

The 2022 issue is not simply that the transformer equation is wrong. The deeper
issue is that the model became too permissive after validation.

The chain is:

\begin{enumerate}[leftmargin=*]
  \item The validation year before 2022 was unusually favorable for broad K2
        exposure.
  \item Because many K2 trades made money in validation, a low threshold looked
        attractive.
  \item That low threshold allowed many trades in 2022.
  \item In 2022, the unconditional T1/K2 environment was poor, so broad
        selection created losses.
\end{enumerate}

This is a risk-control and validation-objective problem as much as a model
architecture problem.

\section{Can Changing the Attention Equation Reduce 2022 Risk?}

Changing the equation can help only if the change gives the model information
that separates dangerous regimes from safe regimes. For example, useful
extensions could include:

\begin{itemize}[leftmargin=*]
  \item market volatility regime distance;
  \item SPY trend or drawdown regime distance;
  \item daily option-liquidity stress distance;
  \item concentration controls by day or expiration;
  \item downside-aware training or threshold selection.
\end{itemize}

But simply making the transformer larger, deeper, or more complex may not solve
the problem. If threshold selection is based only on validation total P\&L, a
single strong validation year can still cause the model to trade too broadly.

For this reason, the next robust version should evaluate and possibly tune
using:

\begin{itemize}[leftmargin=*]
  \item total normalized P\&L;
  \item Sharpe ratio;
  \item Sortino ratio;
  \item maximum drawdown;
  \item worst-year P\&L;
  \item number of selected trades;
  \item maximum trades per day.
\end{itemize}

\section{Sharpe vs Sortino in This Project}

Sharpe ratio is:

\[
  \text{Sharpe}
  =
  \frac{\text{average daily P\&L}}{\text{standard deviation of daily P\&L}}
  \sqrt{252}.
\]

Sortino ratio is:

\[
  \text{Sortino}
  =
  \frac{\text{average daily P\&L}}{\text{standard deviation of negative daily P\&L}}
  \sqrt{252}.
\]

The difference is important. TP2 option strategies can have low win rate but
large positive payoff skew. Sharpe penalizes both upside and downside
volatility. Sortino focuses only on downside volatility.

For this project, Sortino answers:

\begin{center}
\textit{Is the strategy making money while controlling bad downside days, or is
it just relying on a few large winners?}
\end{center}

\section{Simple Example}

Suppose today's TP2 episode is a K2 candidate with:

\begin{itemize}[leftmargin=*]
  \item high TP2 severity;
  \item short D leg near expiration;
  \item acceptable spread;
  \item similar moneyness to several past profitable K2 episodes.
\end{itemize}

The transformer compares it with the previous 64 historical TP2 tokens.

Past episode A is recent, same strategy, similar moneyness, similar liquidity.
It receives high attention.

Past episode B is old, illiquid, and far away in moneyness. Even if its raw
features have some similarity, the multivariate distance penalty lowers its
attention.

Past episode C is T1, not K2, but has the same short-D-leg structure. A
cross-strategy head may still attend to it if that information has historically
helped prediction.

The model then combines these historical memories and outputs:

\[
  \Pr(\text{K2 profitable}) = 0.62.
\]

If the validation-tuned threshold is below 0.62 and T1 is less attractive, the
selector chooses K2. If both T1 and K2 fall below threshold, the episode is
skipped.

\section{Main Takeaways}

\begin{enumerate}[leftmargin=*]
  \item The transformer is a TP2 episode selector, not a TP2 detector.
  \item Each token is an episode-strategy candidate: T1 or K2.
  \item Attention lets the model use similar past TP2 episodes as memory.
  \item MV-ALiBi makes attention financially meaningful by penalizing distance
        in time, option surface, short-D-leg structure, liquidity, and lifecycle.
  \item The 2022 weakness is likely driven more by aggressive validation-time
        selection than by the attention equation alone.
  \item The next research step is not only better architecture, but better
        downside-aware selection using Sortino, drawdown, exposure limits, and
        concentration controls.
\end{enumerate}

\end{document}
